% ---------------------------------------------------------------
\section{Ablation Study}
\label{sec:ablation}


Each ablation removes, replaces or modifies one component and keeps everything else
fixed on the same seeds and streams; the full results, with intervals, are listed in
Table~\ref{tab:ablation} (Appendix~\ref{app:extra}). Fig.~\ref{fig:ablation} plots the two central trade-offs. Panel (a)
shows that every gain in dependence detection costs stationary alarms. Panel (b)
shows that the harm rules form a recall--cost frontier on C6, and that input alarms
used as a harm proxy lie below it. The following paragraphs explain what each
component contributes.

\begin{figure}[t]
\centering
\resizebox{\columnwidth}{!}{% Ablation trade-offs.
% (a) data/ablation_detector.dat <- results/phase3/dependence_extension_20261116_final_v2/{validated_summary,online_dependence_scores}.csv
% (b) data/ablation_rules.dat    <- results/phase3/v2_final/harm_counts.csv (C6, logistic, all batches)
\begin{tikzpicture}
\begin{groupplot}[
  group style={group size=2 by 1, horizontal sep=13mm},
  idm, width=0.52\columnwidth, height=4.8cm, xmajorgrids,
  nodes near coords={\pgfplotspointmeta}, point meta=explicit symbolic,
  visualization depends on={value \thisrow{pos} \as \pos},
  every node near coord/.append style={font=\scriptsize, text=cink, anchor=\pos, inner sep=2.5pt},
]
\nextgroupplot[
  title={(a) Detectors},
  xlabel={stationary alarms / 24 batches}, ylabel={C6 onset detected (\%)},
  xmin=0, xmax=3.2, ymin=35, ymax=115, ytick={40,60,80,100},
]
\draw[-{Stealth[length=2mm]}, cgray, thick, shorten >=5pt, shorten <=5pt] (axis cs:0.8333,53.33) -- (axis cs:2.3667,100);
\addplot[only marks, mark=*, mark size=3pt, c1, fill=c1!70] table[x=x, y=y, meta=label] {data/ablation_detector.dat};
\nextgroupplot[
  title={(b) Harm rules on C6},
  xlabel={false flags (of 306)}, ylabel={harmful batches flagged (\%)},
  xmin=-6, xmax=100, ymin=-10, ymax=112, ytick={0,25,50,75,100},
]
\addplot[only marks, mark=diamond*, mark size=3.2pt, c2, fill=c2!70] table[x=x, y=y, meta=label] {data/ablation_rules.dat};
\end{groupplot}
\end{tikzpicture}
}
\caption{Ablation trade-offs. (a) Detectors on fresh seeds: onset detection of the dependence-only change C6 against the stationary alarm burden (30 seeds each); \emph{Combined} raises an alarm if either detector fires. (b) Harm rules on C6 (logistic monitor, 414 harmful and 306 non-harmful batches): share of harmful batches flagged against false flags. KS-based input alarms used as a harm proxy recover little harm.}
\label{fig:ablation}
\end{figure}

\textbf{Correctness calibration.} Removing the isotonic map changes accuracy estimation
in opposite directions for the two learners. The tree's raw confidence overstates its
correctness, and calibration lowers its pooled MAE from 0.134 to 0.105. Logistic
regression is already close to calibrated on the source law, and the map adds a small
error (0.1142 against 0.1134). Neither version helps under conditional-only change,
because calibration fitted on the source law cannot transfer to a reversed law.

\textbf{Harm rule, estimator and latent tests.} Fig.~\ref{fig:ablation}(b) shows the
recall--cost frontier on C6. Replacing the Poisson-binomial concern score with a point
threshold on $\Aref-\Ahat_t$ raises recall from 0.386 to 0.920 at the cost of 51 rather
than 9 false flags. The ATC-style estimator reaches 0.990 but also raises 23 false
flags on the stationary C0, where the calibrated rule raises none. Under C3--C5 neither the
calibrated nor the point rule fires; the ATC-style rule flags 25 batches in each
condition, of which at most 16 are harmful. Removing the PCA tests halves C6 onset detection (8/30 to 4/30) and leaves every
concern flag unchanged, because the impact channel does not depend on the input alarm
(Algorithm~\ref{alg:loop}).

\textbf{Dependence statistic.} Used on its own, the statistic~\eqref{eq:dep} detects every
C6 onset in a fresh seed set and selects the programmed pair $(x_0,x_1)$ in 30/30 onset windows. On stationary streams it raises 1.70 alarms per 24 batches, against 0.83 for KS/PCA;
its own rate of 70.8 alarms per 1,000 batches exceeds the nominal 50. The development-quantile threshold therefore did
not control the per-batch false-alarm rate. Because the statistic ignores marginal shifts, it cannot replace KS for C1 and C2.
Running both detectors (alarm if either fires) keeps the 30/30 C6 detection, but raises
the stationary burden to 2.37 alarms per 24 batches, 2.8 times that of KS alone.

\textbf{Forecast predictors.} In nested forecasters, monitor features account for most
of the Brier improvement over time alone (0.0644), and released labels add a further
0.0169. On the held-out law, the released-label increment disappears.

\textbf{Adaptation.} Cumulatively refitting the frozen model on released labels does not
change post-onset accuracy under C3 at delay 0, the primary outcome ($+0.02$ pp). In
the late post-onset regime it gains $+2.39$ pp [1.97, 2.87]. A plausible reason is that
each refit mixes the few released reversed labels with the larger pre-reversal training
set. A region-local
trigger appeared to gain 12.45 points over the frozen model. When the update clocks and
label windows were matched to a global refit on the same seeds, however, the local
policy was 0.08 points \emph{worse} than the global one. This result is consistent with
the original global arm having been diluted by pre-reversal training rows; it does not
support a benefit of localization. A
capacity-matched pooled model with cluster interactions reproduced the local gains
almost exactly, including gains on stationary XOR streams. The gain therefore reflects
model capacity rather than drift-aware adaptation.
